\documentclass[10pt,a4paper]{amsart}
\usepackage[margin=19mm]{geometry}
\usepackage{amsmath,amssymb,mathtools}
\usepackage[scr=rsfs,cal=euler]{mathalfa}
\usepackage{xfrac}
\usepackage[unicode,hidelinks,bookmarks=false]{hyperref}
\newcommand{\ie}{i.e.\ }
\newcommand{\cf}{cf.\ }
\newcommand{\resp}{respectively\ }
\newcommand{\Subsection}{\subsection}
\providecommand{\nobreakdash}{\nobreak\hbox{-}}
\setlength{\emergencystretch}{2em}
\multlinegap0pt
\usepackage{bm}
\usepackage{bbm}
\usepackage{dsfont}
\usepackage{xspace}
\usepackage{cancel}
\usepackage{mathtools}
\usepackage{amsthm}
\makeatletter
\@ifundefined{theorem}{\newtheorem{theorem}{Theorem}}{}
\makeatother
\makeatletter
\def\M{{\mathcal M}}
\def\Tr{\operatorname{Tr}}
\newcommand{\brke}[2]{{{\langle #1,#2\rangle}}}
\def\d{{\mathrm d}}
\expandafter\ifx\csname exa\endcsname\relax
\newenvironment{exa}{\begin{example} \rm}{\end{example}}
\fi
\def\parpo#1#2{\{#1,#2\}}
\expandafter\ifx\csname rem\endcsname\relax
\newenvironment{rem}{\begin{remark} \rm}{\end{remark}}
\fi
\makeatother
\title[Poisson quasi-Nijenhuis manifolds, closed Toda lattices, and]{Poisson quasi-Nijenhuis manifolds, closed Toda lattices, and generalized recursion relations}
\author{Eber Chu\~no Vizarreta, Gregorio Falqui, Igor Mencattini, Marco Pedroni}
\date{Source: arXiv:2410.08671v2 (CC BY 4.0)}
\begin{document}
\maketitle

\noindent\emph{Source and licence.} Poisson quasi-Nijenhuis manifolds, closed Toda lattices, and generalized recursion relations. Version arXiv:2410.08671v2 (CC BY 4.0), distributed under the Creative Commons Attribution 4.0 International licence, as recorded in the arXiv licence metadata for this exact version and shown on the abstract page https://arxiv.org/abs/2410.08671v2. Full rights notice: licenses/arxiv-2410.08671-CC-BY-4.0.txt.\par\smallskip

\noindent\textit{Editorial scope.} Counted unit: the Theorem whose source label is \texttt{thm:involution} in the article (source0.tex lines 390--400, source label \texttt{thm:involution}), delivered together with the article's own complete original proof of it (source0.tex lines 401--514, one \texttt{proof} environment of 114 lines). No mathematics is written, repaired or re-derived: statement and proof are the article's own text line for line. Required context retained uncounted: varphi (lines 369--372); recadd (lines 374--378). Every definition, display and label of those lines is retained unchanged. Omitted from this package and not redistributed: 1--368, 373--373, 379--389, 515--1451. Nothing inside a retained excerpt is omitted or altered: every retained body is delivered exactly as the article's lines are written, so no omission receipt of any kind accompanies this package. Citations: the retained excerpts carry no citation command at all, so no bibliography entry of the article is transcribed and no reference is redistributed. Reference adaptation: none was needed and none was made; every reference inside the delivered unit resolves inside it, so no unresolved reference or citation is produced. Printed slips kept literal: no printed word, symbol, hypothesis or number of the counted statement or of its proof was altered. Layout and class: the article's own class is \texttt{article} with packages including amssymb, amsmath, amsthm, calc, xcolor, cancel, authblk, graphicx; the wrapper is the lane's pinned \texttt{amsart} editorial wrapper at 10pt on A4, which restates the theorem-like environments the retained text needs and adds the article's own macro definitions that the retained text needs (\texttt{\textbackslash M}, \texttt{\textbackslash Tr}, \texttt{\textbackslash brke}, \texttt{\textbackslash d}, \texttt{\textbackslash parpo}), with extra wrapper packages (bm, bbm, dsfont, xspace, cancel, mathtools). The delivered numbering is the unit's own. Assets: no retained span contains a figure environment, a graphics-inclusion command, a PDF-page inclusion, a table or a TikZ picture; no image file is copied into this package, the package declares no asset dependency and no image file is redistributed with it. Licence: version arXiv:2410.08671v2 of this article is distributed under the Creative Commons Attribution 4.0 International licence, as recorded in the arXiv licence metadata for this exact version and shown on the abstract page https://arxiv.org/abs/2410.08671v2. A full rights notice is delivered at licenses/arxiv-2410.08671-CC-BY-4.0.txt. Characters: every retained excerpt is pure ASCII, so the delivered engine, XeLaTeX with its default Latin Modern text font, reports no missing character.
\begin{equation}
\label{varphi}
\brke{\phi_k}{ X}= \frac12\Tr\left((i_XT_N)\, N^{k}\right)= \frac12\Tr\left(N^{k}\, (i_XT_N)\right),\qquad k\ge 0.
\end{equation}

\begin{equation}
\label{recadd}
\parpo{H_k}{H_j}-\parpo{H_{k-1}}{H_{j+1}}=-\brke{ \phi_{j-1}}{\pi\, \d H_{k-1}}-\brke{\phi_{k-2}}{ \pi\, \d H_j},
\qquad k>j \ge 1,
\end{equation}

\begin{theorem} 
\label{thm:involution}
Let $(\M,\pi,N,\phi)$ be a PqN manifold,  
and $H_k=\frac1{2k}\Tr({N}^k)$. Suppose that there exists a 2-form $\Omega$ such that:
\begin{itemize}
\item[$(a)$] $\phi=-2\,\d H_1\wedge\Omega$; 
\item[$(b)$] $\Omega(X_j,{Y_k})=0$ for all $j,k\ge 1$, where $Y_k=N^{k-1}X_1-X_k$ and $X_k=\pi^\sharp\,\d H_k$.
\end{itemize}
Then %the PqN manifold is involutive, i.e., 
$\parpo{H_j}{H_k}=0$ for all $j,k\ge 1$. 
\end{theorem}

\begin{proof} Since on a PqN manifold, for all vector fields $X,Y$, 
\begin{equation}
\label{eq:pre-ThatN}
T_{N}(X,Y)=\pi^\sharp(i_{X\wedge Y}\phi),
\end{equation}
we have that
\begin{equation}
\label{eq:ThatN}
\begin{split}
T_{N}(X,Y)& %=\pi^\sharp(i_{X\wedge Y}d_N\Omega)
\stackrel{(a)}{=}-2\pi^\sharp(i_Yi_X (\d H_1\wedge\Omega))
=-2\pi^\sharp(i_Y(\langle \d H_1,X\rangle\Omega-\d H_1\wedge i_X\Omega))\\
&=-2\pi^\sharp(\langle \d H_1,X\rangle i_Y\Omega-\langle \d H_1,Y\rangle i_X\Omega+(i_Yi_X\Omega) \d H_1)\\
&=-2\langle \d H_1,X\rangle(\pi^\sharp\,\Omega^\flat)(Y)+2\langle \d H_1,Y\rangle(\pi^\sharp\,\Omega^\flat)(X)-2\Omega(X,Y)X_1,
\end{split}
\end{equation}
%for all vector fields $X,Y$, 
so that 
\begin{equation}
\label{eq:5}
i_XT_{N}=-2\langle \d H_1,X\rangle\,\pi^\sharp\,\Omega^\flat
+2(\pi^\sharp\,\Omega^\flat)(X)\otimes \d H_1-2X_1\otimes \Omega^\flat X.
\end{equation}
Therefore %, by the definition %(\ref{varphi}) 
%of the 1-forms $\phi_k$, 
\begin{equation}
\label{eq-pre-tr}
\begin{split}
\brke{\phi_k}{X} &\stackrel{\eqref{varphi}}{=}\frac12\Tr\left({N}^k(i_{X}T_{N})\right)
=%\frac12
\Tr\left({N}^k\left(-\langle \d H_1,X\rangle\,\pi^\sharp\,\Omega^\flat
+(\pi^\sharp\,\Omega^\flat)  (X)\otimes \d H_1
-X_1\otimes \Omega^\flat X\right)\right)\\
&=-%\frac12
\langle \d H_1,X\rangle\Tr\left({N}^k\pi^\sharp\,\Omega^\flat\right)
+%\frac12
\Tr\left(({N}^k\pi^\sharp\,\Omega^\flat)(X)\otimes \d H_1\right)
-%\frac12
\Tr\left(({N}^k X_1)\otimes \Omega^\flat X\right).
\end{split}
\end{equation}
Now observe that the last two terms in the last line of \eqref{eq-pre-tr} sum up to $-2\Omega(X,N^kX_1)$. In fact,
\begin{equation}
\Tr(X\otimes\alpha)=\brke{\alpha}{X}\label{eq:identi} 
\end{equation} 
entails that $-%\frac12
\Tr\left(({N}^k X_1)\otimes \Omega^\flat X\right)=-\Omega(X,N^kX_1)$ and
\begin{equation}
\label{eqtr-b}
\begin{split}
\Tr\left(({N}^k\pi^\sharp\,
\Omega^\flat)(X)\otimes \d H_1\right)
&=\brke{\d H_1}{({N}^k\pi^\sharp\,
\Omega^\flat)(X)}
=\brke{\d H_1}{(\pi^\sharp ({N}^*)^k \Omega^\flat)(X)}\\
&=-\brke{(({N}^*)^k \Omega^\flat)(X)}{X_1}
=-\brke{\Omega^\flat(X)}{{N}^k X_1}
=-\Omega(X,{N}^k X_1).
\end{split}
\end{equation}
The previous remarks imply that, for every vector field $X$,
\begin{equation}
\label{eqtr}
\brke{\phi_k}{X}=%\Tr\left({N}^k(i_{X}T_{N})\right)=
-\langle \d H_1,X\rangle\Tr\left({N}^k\pi^\sharp\,\Omega^\flat\right)-2\Omega(X,{N}^k X_1),
\end{equation}
meaning that
\begin{equation}
\label{eq-phi_k}
\phi_k=2\Omega^\flat({N}^k X_1)-\Tr\left({N}^k\pi^\sharp\,\Omega^\flat\right)\d H_1.
\end{equation}
Putting %instead 
$X=X_j$ in (\ref{eqtr}), one has
\begin{equation}
\label{eqtr2}
\brke{\phi_k}{X_j}=-\parpo{H_1}{H_j}\Tr\left({N}^k\pi^\sharp\,\Omega^\flat\right)-2\Omega(X_j,{N}^k X_1).
\end{equation}
To prove that the traces $H_k$ of the powers of $N$ are in involution, we show by induction on $n$ 
%[ovviamente $n$ non e' il numero delle particelle, il teorema non riguarda (solo) Toda] 
that the following property holds:
\begin{equation*}
%\label{inv-induzione}
(P_n):\qquad\parpo{H_l}{H_m}=0\quad\mbox{for all pairs $(l,m)$ such that $l+m\le n$.}
\end{equation*}
This is obviously true for $n=2$. Let us show that $(P_n)$ implies $(P_{n+1})$. If $l+m=n+1$ and $l>m$, then 
\begin{equation}
\label{recadd2}
\begin{split}
\parpo{H_l}{H_m}&\stackrel{\eqref{recadd}}{=}\parpo{H_{l-1}}{H_{m+1}}-\brke{ \phi_{m-1}}{X_{l-1}}-\brke{\phi_{l-2}}{X_m}\\
&\stackrel{\eqref{eqtr2}}{=}\parpo{H_{l-1}}{H_{m+1}}+\parpo{H_1}{H_{l-1}}\Tr\left({N}^{m-1}\pi^\sharp\,\Omega^\flat\right)+2\Omega(X_{l-1},{N}^{m-1} X_1)\\
&\ \ \ +\parpo{H_1}{H_m}\Tr\left({N}^{l-2}\pi^\sharp\,\Omega^\flat\right)+2\Omega(X_m,{N}^{l-2} X_1).
\end{split}
\end{equation}
Since $\parpo{H_1}{H_{l-1}}=\parpo{H_1}{H_m}=0$ by the induction hypothesis, we obtain
\begin{equation}
\label{recadd3}
%\begin{split}
\parpo{H_l}{H_m}%&
=\parpo{H_{l-1}}{H_{m+1}}+2\Omega(X_{l-1},{N}^{m-1} X_1)+2\Omega(X_m,{N}^{l-2} X_1).
%\end{split}
\end{equation}
Now, thanks to assumption (b), we can substitute ${N}^{i-1}X_1$ with $X_{i}$ in the 
last two terms, showing that their sum vanishes. Hence we obtain that 
\begin{equation}
\label{recadd4}
\parpo{H_l}{H_m}=\parpo{H_{l-1}}{H_{m+1}}.
\end{equation}
In the same way, we can show that $\parpo{H_{l-1}}{H_{m+1}}=\parpo{H_{l-2}}{H_{m+2}}$ and so on, until we obtain that 
\begin{equation}
\label{recadd5}
\parpo{H_l}{H_m}=\parpo{H_{m}}{H_{l}}=0,
\end{equation}
so that $(P_{n+1})$ holds.
\end{proof}  %\hfill$\square$ \medskip

\end{document}
